% !TeX root = ../教材题库_全部条目_转写工作稿.tex
% Source: XJ-B1-C04-EX-P115-Y1132
% 第4章 / 4.3 对数函数 / 教材 p.115、116 / PDF 第122、123页
% OCR未校对：公式与文字均可能有误；未完成数学TeX转写；禁止直接发给学生
% 原题及相邻原解截图（核对用）：../题目原图/XJ-B1-C04-EX-P115-Y1132-01.webp, ../题目原图/XJ-B1-C04-EX-P115-Y1132-02.webp
% 本文件只在转写完成后作为题干使用；OCR 条目尚不是可靠题干。
\begin{bankitem}{4.3 例题 例题1}{第4章 / 4.3 对数函数 / 教材 p.115、116 / PDF 第122、123页}{XJ-B1-C04-EX-P115-Y1132}{OCR未校对：公式与文字均可能有误；未完成数学TeX转写；禁止直接发给学生}
\begin{ocrdraft}
[教材p.115 / PDF第122页]
\par
00将下列指数式化为对数式，对数式化为指数式：
\par
23一8；
\par
105=100000；
\par
（1）
\par
（2）
\par
（3）
\par
3上
\par
（4）
\par
log232一5；
\par
1
\par
（5）
\par
log3
\par
（6）
\par
Iogxb=2．
\par
27
\par
解
\par
（Dlog28一3；
\par
（2）
\par
log10100000=5；
\par
（3）
\par
10g37一工；
\par
（4）
\par
25一32；
\par
1
\par
（5）
\par
3一
\par
（6）=乃（>0且乒1).
\par
27，
\par
把对数定义中的b=lwN代人'一N，得到
\par
alogaN==N（N>0，a>0且乒1)；
\par
把N=ab代人b=logN,得到
\par
[教材p.116 / PDF第123页]
\par
乃一lo(b\ensuremath{\text{[OCR:EUR]}}-R,>0且1）．
\par
上述两个等式叫作对数的基本恒等式．用它们做有关对数的推理或计算比较方
\par
便．例如，由基本恒等式立刻知道：
\par
logua=log\ocrglyph{201E}a
\par
log\ocrglyph{201E}l=log\ocrglyph{201E}a
\par
0
\par
即底的对数为1，1的对数为0．
\end{ocrdraft}
\end{bankitem}
